\section{依存分析的评估}

\subsection{UAS 与 LAS}

\begin{figure}[H]
\centering
\includegraphics[width=0.95\textwidth]{slides-images/slide-031.jpg}
\caption{依存分析评估示例：UAS 和 LAS\protect\footnotemark}
\end{figure}
\footnotetext{Slides 第31页。}

依存分析的评估非常直观——逐条检查系统预测的依存弧是否与标准答案一致：

\begin{importantbox}{两种评估指标}
\begin{itemize}
    \item \textbf{UAS（Unlabeled Attachment Score）}：只看依存弧的方向是否正确（即每个词的中心词是否预测对了），不考虑关系标签
    $$\text{UAS} = \frac{\text{正确预测中心词的词数}}{\text{总词数}}$$
    \item \textbf{LAS（Labeled Attachment Score）}：同时要求依存弧方向和关系标签都正确
    $$\text{LAS} = \frac{\text{中心词和标签都正确的词数}}{\text{总词数}}$$
\end{itemize}
LAS $\leq$ UAS，因为 LAS 的要求更严格。
\end{importantbox}

例如对句子 ``She saw the video lecture''（5 个词），如果系统有 4 个词的中心词预测正确（1 个错误），但只有 2 个词的标签也完全正确，则 UAS = 4/5 = 80\%，LAS = 2/5 = 40\%。

\subsection{本章小结}

UAS 和 LAS 是依存分析的标准评估指标。UAS 只评估依存弧方向的正确性，LAS 同时要求方向和标签都正确。这种基于树库的定量评估方法是 1990 年代以来 NLP 领域的重要进步。

%% ========== 后续发展 ==========

